\chapter{Literature Review}
\label{ch:litreview}

This chapter is organised around the conceptual problem the dissertation addresses, rather than as a list of individually summarised papers. \S\ref{sec:lit-representation} discusses how representation is treated in customer segmentation research. \S\ref{sec:lit-validation} discusses cluster structure and validation. \S\ref{sec:lit-explainable} distinguishes three design positions in explainable clustering. \S\ref{sec:lit-complexity} discusses how those methods measure and control explanation complexity. \S\ref{sec:lit-gap} states the resulting research gap and positions this dissertation within it.

\section{Representation in customer segmentation}
\label{sec:lit-representation}

Research on customer segmentation often treats representation as a \emph{neutral input} to clustering — that is, it implicitly assumes that the particular choice of which behavioural signals to encode is a preliminary, largely inconsequential step, and that any reasonable encoding of the same customers would support broadly the same substantive conclusions. That assumption is difficult to defend for \emph{event logs} — timestamped records of individual visitor actions such as page views, cart additions and transactions, of the kind analysed in this dissertation (see \S\ref{sec:data-source}). A customer can be described by a handful of purchase summaries, by the full raw sequence of recorded actions, or by a richer set of temporal, conversion and category signals computed from that sequence; each choice changes the distances the clustering algorithm observes between the same people. Clickstream research describes raw action sequences as large, variable and noisy, and instead frames useful segments as subsets recovered through constraints on \emph{derived} attributes \citep{dextras2019segmentifier}. Work on object-centred event logs likewise converts raw traces into vector profiles before any clustering or pattern analysis is attempted \citep{faria2023eventlogs}. Customer behaviour studies construct features for segmentation and subsequent decision tasks, but their application-specific choices do not by themselves establish a universally correct representation \citep{abdi2019cbmf}. Representation is therefore better understood as part of the empirical model being fitted to customer event logs, rather than as a preliminary clerical step that precedes the "real" modelling decisions.

\section{Cluster structure and validation}
\label{sec:lit-validation}

Clustering does not reveal a representation-independent truth about a population. Clusterability theory supplies several formally incompatible notions of when data possess recoverable structure at all \citep{ackerman2009clusterability}, and empirical comparisons show that clusterability measures themselves depend on distributional conditions, dimensionality and the type of structure sought \citep{adolfsson2019clusterability}. A high internal validity score therefore supports a partition \emph{within} the feature space it was computed in — it is not evidence of an intrinsic segmentation that exists independently of that space. Dimensionality reduction retains this dependence rather than removing it: it replaces named variables with a new, rotated geometry, and in doing so can obscure the link between a resulting cluster and the original, business-meaningful attributes \citep{bertsimas2021icot}. A paired comparison on the same customers — the design adopted in this dissertation — is accordingly more informative about the effect of representation than comparing results drawn from unrelated populations or studies, because it holds the population fixed and varies only the representation.

\section{Explainable clustering}
\label{sec:lit-explainable}

Explainable clustering methods are often discussed as a single category, but they occupy at least three distinct design positions, which must not be conflated.

\textbf{Intrinsic} methods expose the computation that produces a partition, as part of producing it. CLASSIX sorts observations along a principal direction (the direction of maximum variance in the standardised data), aggregates nearby points into groups starting from a sequence of starting points, and merges the resulting groups according to a stated rule \citep{chen2024classix}; its public group assignments, starting points and final labels together support a step-by-step trace from any one observation to its final cluster. \textbf{Explainability-first} methods instead build a simple description directly into the clustering objective, so that simplicity is optimised for jointly with cluster quality rather than added afterwards. ICOT constructs tree-shaped partitions through combinatorial optimisation, while XClusters balances the usual clustering distortion objective against the size and accuracy of an accompanying tree during the formation of the clusters themselves \citep{bertsimas2021icot,hwang2023xclusters}. \textbf{Post hoc} methods begin from an already-fitted reference partition and approximate it afterwards with rules or a tree, without influencing how that reference partition was formed; ExKMC, used later in this dissertation, is one such method.

These three positions answer different questions about different objects, and their outputs are not interchangeable. Provenance from an intrinsic algorithm answers how its \emph{own} grouping arose. A post hoc tree answers how closely a set of axis-aligned rules can reproduce \emph{another}, already-fixed partition. Neither answer can be converted into the other without changing what is being explained — a point developed further for CLASSIX and ExKMC specifically in \S\ref{sec:lit-gap} and in Chapter~\ref{ch:results}.

{\setlength{\tabcolsep}{4pt}
\begin{longtable}{p{0.12\textwidth}p{0.14\textwidth}p{0.14\textwidth}p{0.15\textwidth}p{0.17\textwidth}p{0.14\textwidth}}
\caption{Comparison of explanation designs in the reviewed literature. \emph{Timing} records whether the explanation is produced intrinsically as clustering proceeds, jointly with the clustering objective, or post hoc after clustering is already complete. \emph{Complexity control} records the specific setting or budget through which each method limits how elaborate its explanation can become.}\label{tab:literature_designs}\\
\toprule
Method & Target & Timing & Form & Complexity control & Status \\
\midrule
\endfirsthead
\toprule
Method & Target & Timing & Form & Complexity control & Status \\
\midrule
\endhead
CLASSIX & Own partition & Intrinsic & Groups and starting points & Radius and minimum size & Journal \\
IMM & K Means objective & Tree construction & Threshold tree & One leaf per cluster & Conference \\
ExKMC & K Means labels & Post hoc & Expanded tree & Leaf budget & Preprint \\
CREAM & Fitted partition & Post hoc & Cubes and tree rules & Hypercube tuning parameters (coverage/purity thresholds) & Conference \\
ICOT & Own partition & Joint & Optimised tree & Validation objective and tree size & Journal \\
XClusters & Joint objective & Joint & Decision tree & Distortion and node tradeoff & Conference \\
\bottomrule
\end{longtable}
}

\section{Explanation complexity and fidelity}
\label{sec:lit-complexity}

Decision trees make the tradeoff between concise expression and partition quality explicit. IMM uses threshold cuts to construct a small tree and analyses the additional K Means or K Medians cost that restriction causes \citep{moshkovitz2020imm}. Theoretical work shows that this kind of simplicity can impose an unavoidable loss in clustering quality, and that the size of that loss depends on the specific clustering problem, the number of clusters and the dimensionality \citep{laber2021price}. Related work asks whether a fixed clustering can be made explainable after removing a limited number of outlying observations, which further demonstrates that an apparently simple tree can obtain its simplicity by quietly changing the population being explained, rather than by finding a genuinely simpler description of the original population \citep{bandyapadhyay2022explanation}. CREAM constructs hypercubic approximations through decision-tree induction and tunes the resulting approximation directly, extending the same concern to a second reference algorithm \citep{sabbatini2023cream}. Together, these methods make complexity measurable — through leaf count, tree depth, the number of conditions per rule, and the number of features used — but that measurability does not by itself establish that a person will understand, trust or act on the resulting explanation.

Theory and software maturity vary across these methods in ways that matter for how their results can be used. IMM and the price-of-explainability analyses above provide mathematical guarantees on how much clustering quality can be lost by imposing a simple tree, under precisely defined objective functions. CLASSIX, by contrast, supplies a maintained software implementation and public explanation attributes for its own fitted partition, but it does not provide the same kind of quality-loss guarantee. CREAM, ICOT and XClusters each broaden the space of allowed explanations or add joint optimisation of size and accuracy, but their formal guarantees and their computational cost concern different targets from either CLASSIX or IMM. The practical implication is that queryable rules, a bounded loss in clustering quality, and native intrinsic provenance are three separate, complementary properties of an explanation method — not three points on one common "explainability" scale that would let these methods be ranked against each other.

ExKMC expands a decision tree beyond one leaf per K Means cluster, so that the analyst can test whether allowing more terminal nodes (a larger \emph{leaf budget}) improves agreement with a fixed reference partition. Its only verified source is an arXiv preprint rather than a peer-reviewed publication; it documents the implemented method but does not itself provide independent evidence of general effectiveness \citep{frost2020exkmc}. \emph{Fidelity}, as used throughout this dissertation, means the proportion of exact label agreements between the tree's predictions and a set of held-out K Means assignments — it is neither an internal clustering quality metric nor evidence that the underlying K Means partition is substantively correct, only a measure of how closely the tree reproduces it. More general explainable-AI research finds that fidelity measures can disagree with each other even when a fully transparent reference model is available, so the specific target being approximated and the specific calculation used must always be stated explicitly rather than assumed \citep{mironicolau2025fidelity}.

Operational usefulness introduces a further, distinct boundary. A decision tree or a set of threshold conditions can be directly queried in a database, while an intrinsic trace can instead identify representative points and merge relationships. \emph{Queryability} can be observed through the number of conditions in a rule, whether those conditions use named (business-meaningful) features, how many customers the rule covers, and how closely it agrees with a reference partition — but none of these observable properties by itself establishes improved revenue, a more effective intervention, or genuine human comprehension. Segmentifier illustrates why recorded constraints and provenance can support the interactive refinement and export of behavioural segments, but its design study does not attempt to validate the commercial effect of the resulting clustering model \citep{dextras2019segmentifier}. Equally, a theoretical guarantee on objective cost does not by itself measure whether a resulting segment is large enough, sufficiently distinct, or covered by reliable enough attributes to be acted upon in practice — this last question is addressed directly, for the RetailRocket clusters, by the addressability audit introduced in \S\ref{sec:candidates} and reported in \S\ref{sec:explanation-results}.

\section{Research gap and positioning of this dissertation}
\label{sec:lit-gap}

Two gaps follow from the discussion above, and both motivate the design of this dissertation. First, the effect of representation on clustering is rarely isolated directly: it is usually varied incidentally alongside algorithm choice, cohort, or validation rule, rather than compared for the same customers under matched grids and constraints. Second, and specifically for the two explanation methods used later in this dissertation, CLASSIX provenance and ExKMC fidelity answer genuinely different questions and must not be interpreted as one common explainability score: CLASSIX provenance is exact evidence about how CLASSIX's \emph{own} fitted partition was produced, whereas ExKMC fidelity is evidence about how well a tree can imitate a \emph{separately fitted} K Means partition. A high ExKMC fidelity says nothing about how faithfully CLASSIX's aggregation groups trace their own clusters, and a large or small number of CLASSIX aggregation groups says nothing about how simple a post hoc tree for K Means on the same data would need to be. In line with these observations, the present study compares compact transaction and rich behavioural representations for the same RetailRocket purchasers under matched clustering and evaluation conditions, and relates the resulting differences in cluster structure to CLASSIX provenance, ExKMC fidelity and auditable addressability as three separate, non-interchangeable strands of evidence.
